\chapter*{Annexe B -- Manifests Kubernetes}
\addcontentsline{toc}{chapter}{Annexe B -- Manifests Kubernetes}
\label{annexe-k8s}

Les manifestes ci-dessous correspondent aux ressources Kubernetes réellement déployées sur le cluster k3s du VPS OVH, référencées au \cref{chap-sprint4}. \emph{Par mesure de sécurité, les valeurs du manifeste \texttt{Secret} sont remplacées par des espaces réservés~: les valeurs réelles ne doivent jamais figurer dans ce rapport ni dans le dépôt Git du projet.}

\begin{lstlisting}[language=bash,caption={Deployment et Service -- Backend}]
apiVersion: apps/v1
kind: Deployment
metadata:
  name: backend
  namespace: monitoring
spec:
  replicas: 1
  selector:
    matchLabels:
      app: backend
  template:
    metadata:
      labels:
        app: backend
    spec:
      containers:
      - name: backend
        image: mariemchaabani/monitoring-backend:<tag>
        ports:
        - containerPort: 3000
        envFrom:
        - secretRef:
            name: monitoring-secret
        volumeMounts:
        - name: backups
          mountPath: /app/backups
      volumes:
      - name: backups
        hostPath:
          path: /home/debian/DevOps-Monitoring-App/backups
---
apiVersion: v1
kind: Service
metadata:
  name: backend
  namespace: monitoring
spec:
  selector:
    app: backend
  ports:
  - port: 3000
    targetPort: 3000
  nodePort: 30300
  type: NodePort
\end{lstlisting}

\begin{lstlisting}[language=bash,caption={Deployment et Service -- Frontend}]
apiVersion: apps/v1
kind: Deployment
metadata:
  name: frontend
  namespace: monitoring
spec:
  replicas: 1
  selector:
    matchLabels:
      app: frontend
  template:
    metadata:
      labels:
        app: frontend
    spec:
      containers:
      - name: frontend
        image: mariemchaabani/monitoring-frontend:<tag>
        ports:
        - containerPort: 80
---
apiVersion: v1
kind: Service
metadata:
  name: frontend
  namespace: monitoring
spec:
  selector:
    app: frontend
  ports:
  - port: 80
    targetPort: 80
    nodePort: 30080
  type: NodePort
\end{lstlisting}

\begin{lstlisting}[language=bash,caption={Secret -- configuration sensible (valeurs redigees)}]
apiVersion: v1
kind: Secret
metadata:
  name: monitoring-secret
  namespace: monitoring
type: Opaque
stringData:
  MONGODB_URI: "<mongodb-atlas-connection-string>"
  EMAIL_USER: "<email-expediteur>"
  EMAIL_PASS: "<mot-de-passe-application>"
  PORT: "3000"
  JWT_SECRET: "<cle-secrete-jwt>"
  ADMIN_EMAIL: "<email-admin>"
  ADMIN_PASSWORD: "<mot-de-passe-admin>"
\end{lstlisting}
